\documentclass[11pt]{article}
\usepackage[margin=1in]{geometry}
\usepackage{amsmath,amssymb}
\usepackage{hyperref}
\usepackage{longtable}
\usepackage{tikz}
\usetikzlibrary{arrows.meta,positioning}

\newcommand{\code}[1]{\texttt{#1}}
\newcommand{\notimpl}{\textbf{Not implemented.}}
\newcommand{\notest}{\textit{Not established --- see \code{docs/OPEN\_QUESTIONS.md}.}}

\title{PINTHAC Code Summary}
\author{Elliott Hirko}
\date{Phase 8 revision}

\begin{document}
\maketitle

\begin{abstract}
This is the model manual for PINTHAC (\emph{Physics Informed Nuclear
Thermal-Hydraulic Analysis Code}). It keeps the section numbering of the original
project table of contents, \code{docs/reference/PINTHA\_Code\_Summary.pdf}, and fills
every section from the docstrings that now live in \code{pinthac/}. Nothing in this
document was re-derived: every formula, valid range, uncertainty figure and reference
below is transcribed from the corresponding function's docstring in the code, which is
itself traced to \code{docs/reference/} or reported as \notest{} where no source
document supplies it. A section left blank in the original PDF and still blank here
means the model is not implemented; that is stated explicitly rather than silently
omitted. See \code{docs/FINAL\_REPORT.md} for what changed, phase by phase, and
\code{docs/OPEN\_QUESTIONS.md} for everything still open.

Sections and functions the PDF's table of contents did not anticipate, but that exist
in the code and are worth documenting, are added as lettered sub-subsections (e.g.
Zircaloy under Steels, Mikityuk under Liquid Sodium) rather than renumbering anything
the PDF already had.
\end{abstract}

\tableofcontents

% =====================================================================================
\section{Properties}
% =====================================================================================

Material and fluid property models. Every function is a flat-namespace class member
(\code{UO2.k\_Klimenko(T)}, no \code{self}), accepting a Python float, a NumPy array, or
a PyTorch tensor and returning the same type, per \code{CLAUDE.md} section 5.

\subsection{Water Properties}

Implemented in \code{pinthac/properties/iapws95.py} (the 1995 equation of state) and
\code{pinthac/properties/iapws97.py} (the 1997 industrial formulation, and the
viscosity/thermal-conductivity release formulations \code{iapws95.py} calls into).
Every function in these two modules that touches core Helmholtz-energy arithmetic has a
thin docstring (a one-line description, no formulation/range/uncertainty block) --- the
physics is instead verified directly against 41 published IAPWS check values in
\code{tests/test\_iapws\_verification.py}, which is the authoritative statement of what
is and is not correct in this implementation. This section summarizes that verification
rather than re-deriving the (very long) Helmholtz free-energy expansions themselves.

\subsubsection{IAPWS 1995 Equation of State}

\code{properties/iapws95.py::IAPWS95}. The full Helmholtz free-energy formulation
(ideal-gas part \code{Phi0} plus residual part \code{Phir}), from which pressure,
entropy, enthalpy, isochoric and isobaric heat capacity, and speed of sound are all
analytic derivatives (\code{p}, \code{s}, \code{h}, \code{cv}, \code{cp}, \code{c}).
Built on PyTorch throughout, so every one of those derivatives is available via
autograd, not finite differences, and every function stays differentiable under a
batch of tensor inputs.

Two solves sit on top of the analytic core, both implemented as fixed-iteration-count,
batch-vectorized loops (no \code{scipy.optimize}, per \code{CLAUDE.md} section 5.5):
\begin{itemize}
  \item \code{saturation(T)}: the coexisting liquid/vapor state at a given temperature,
    via Newton's method on the 2$\times$2 Maxwell-construction system (equal pressure,
    equal Gibbs energy), seeded by the IAPWS-95 ancillary equations and using an
    autograd-built Jacobian.
  \item \code{rho\_Tp(T, p)}: density at given temperature and pressure, via damped
    Newton iteration seeded by a physically motivated initial guess (the saturated-liquid
    ancillary density for compressed liquid, ideal-gas density otherwise). The
    docstring is explicit that a global bisection search over the full density range is
    \emph{not} safe here: away from the true $(T,\rho)$ branch the residual terms (some
    with $\tau$ exponents up to 50) stop cancelling cleanly in floating point, and a
    bracket search can lock onto a spurious root near the critical density instead of
    the physical one.
  \item \code{T\_hp(h, p)}: temperature at given enthalpy and pressure, via bisection on
    \code{rho\_Tp}+\code{h}, valid where enthalpy is monotonic in temperature at fixed
    pressure (true on the supercritical isobars used throughout the SCA solvers; not
    valid inside the two-phase dome).
\end{itemize}

\textbf{Valid range / Uncertainty:} not stated per-function in the code; established
instead by the verification suite below.

\textbf{Reference:} IAPWS R6-95(2018), \emph{Revised Release on the IAPWS Formulation
1995 for the Thermodynamic Properties of Ordinary Water Substance for General and
Scientific Use}.

\textbf{Verification (\code{tests/test\_iapws\_verification.py}):} all of R6-95's Table
7 (single-phase region: $p$, $c_v$, speed of sound, entropy, 11 points spanning
300--900~K and 0.1--700~MPa) and Table 8 (saturation: $p_{sat}$, $\rho_f$, $\rho_g$,
$h$, $s$) reproduce the release's own values to its stated nine significant figures.
This is the core claim of the GPU IAPWS-95 project and it now has an external proof.
One conditioning limit, not a defect: at $T=647$~K, $\rho=358$~kg/m$^3$ (0.6~K below
$T_c$, near $\rho_c$), $c_v$ and $s$ remain exact to $5\times10^{-10}$ and
$8\times10^{-10}$ while pressure is off by $5.3\times10^{-6}$ and speed of sound by
$4.6\times10^{-4}$ --- first-order quantities exact, second-order quantities degraded,
the signature of numerical ill-conditioning near the critical point rather than a wrong
formulation.

\subsubsection{IAPWS 1997 Equation of State}

\code{properties/iapws97.py}. \textbf{Region 1} (subcooled liquid, class \code{R1}) and
\textbf{Region 2} (superheated steam, class \code{R2}) are implemented, each with the
forward Gibbs-energy formulation and the backward equations $T(p,h)$ and $T(p,s)$
(Region 2's backward equations auto-select sub-region 2a/2b/2c). The
saturation-pressure/temperature relation (class \code{RSAT}) is also implemented.
\textbf{Region 3} (near-critical) and \textbf{Region 5} (high-temperature steam,
$>$1073~K) are \notimpl{} No formulation, range or uncertainty is stated in the
docstrings for these regions; IAPWS-95 (above) is the differentiable, full-range
equation of state this library actually uses for supercritical-water work, so Regions 1
and 2 here serve mainly the industrial-formulation call sites that still use them (none
of the SCA solvers do).

\textbf{Reference:} IAPWS R7-97(2012), \emph{Revised Release on the IAPWS Industrial
Formulation 1997 for the Thermodynamic Properties of Water and Steam}.

\subsubsection{IAPWS Thermal Conductivity and Viscosity}

\code{properties/iapws97.py::VISC} (R12-08, viscosity) and \code{::COND} (R15-11,
thermal conductivity), both reached from \code{IAPWS95.mu}/\code{IAPWS95.lam}. Each
formulation is a sum of a dilute-gas term, a finite-density (``background'') term, and
a critical-enhancement term active only near the critical point:

\begin{align*}
  \mu(\rho,T) &= \mu_0(T)\cdot\mu_1(\rho,T)\cdot\mu_2(\rho,T) &\text{(R12-08 Eq. 10)}\\
  \lambda(\rho,T) &= \lambda_0(T) + \lambda_1(\rho,T) + \lambda_2(\rho,T,\ldots) &\text{(R15-11)}
\end{align*}

The viscosity's critical enhancement $\mu_2$ needs the isothermal compressibility
$(\partial\rho/\partial p)_T$ at both the state's own temperature and a fixed reference
temperature $T_R=1.5\,T_c=970.644$~K; both come from IAPWS-95 itself (computed on the
\code{iapws95.py} side and passed down), which roughly doubles the cost of a viscosity
evaluation with the enhancement on. \code{IAPWS95.mu(d, enhancement=False)} gives the
industrial simplification $\mu_2=1$ that R12-08 sanctions outside the near-critical
region. The enhancement is left \emph{on} by default because the thermal-conductivity
critical enhancement divides by $\mu$ (R15-11 Eq. 18), so a $\mu$ missing its own
enhancement silently inflates $\lambda$ near the critical point.

\textbf{Verification:} R12-08 Table 4 (all 11 points, away from the critical region) to
$10^{-8}$; R15-11's dilute-gas and dense-liquid end-points of the 647.35~K isotherm.
\textbf{Three confirmed defects, all confined to the critical-enhancement machinery}
(full detail and error tables in \code{docs/OPEN\_QUESTIONS.md}, Round 6):
\begin{enumerate}
  \item Viscosity $\mu_2$ is under-computed near $\rho_c$: error ranges from
    $-0.0003\%$ at $\rho=122$~kg/m$^3$ to $-8.42\%$ at $\rho=322$~kg/m$^3$ (the critical
    density itself), confined to R12-08 Eqs.\ (14)--(19).
  \item Thermal conductivity $\lambda_2$ is over-computed near $\rho_c$ (the opposite
    sign) and returns \textbf{NaN exactly at $\rho_c=322$}, confined to R15-11
    Eqs.\ (17)--(25).
  \item $\lambda_2$ is not zeroed where $\Delta\chi<0$ as R15-11 Sec.\ 2.7 requires:
    $-0.094\%$ at 998~kg/m$^3$ and $-0.190\%$ at 1200~kg/m$^3$, both at 298.15~K.
\end{enumerate}
\textbf{Consequence:} the supercritical-water work in this library runs at 25~MPa,
above $p_c$, so it never sits at the critical point itself, but the pseudocritical
region at 25~MPa is close enough that $\mu$ and $\lambda$ carry some of this error, and
the NaN is reachable. Nothing downstream should be trusted within a few kg/m$^3$ of the
critical density until these are fixed. This is flagged, not fixed, here --- fixing it
touches \code{pinthac/}, out of scope for this phase.

\textbf{Reference:} IAPWS R12-08, \emph{Release on the IAPWS Formulation 2008 for the
Viscosity of Ordinary Water Substance}; IAPWS R15-11, \emph{Release on the IAPWS
Formulation 2011 for the Thermal Conductivity of Ordinary Water Substance}.

\subsection{Liquid Metal Properties}

\code{properties/liqprops.py::Sodium}, \code{::Lead}, \code{::LBE}. Each carries six
properties (density, surface tension, specific heat, specific enthalpy, viscosity,
thermal conductivity), each with its own validated temperature range
(\code{range\_rho}, \code{range\_cp}, \ldots) and relative model-form uncertainty band
(\code{uncert\_rho}, \code{uncert\_cp}, \ldots, each a \texttt{[low, high]} pair) as
class attributes next to the correlation. \code{Props(mat, T)} bundles all six into the
uniform dict \code{getprop.py} dispatches to; \code{RePr(G, D, prop)} computes Reynolds
and Prandtl numbers from that dict.

\begin{longtable}{p{2.1cm}p{4.4cm}p{2.3cm}p{2.3cm}p{3.5cm}}
\textbf{Property} & \textbf{Formulation} & \textbf{Valid range} & \textbf{Uncertainty} & \textbf{Notes} \\ \hline
$\rho(T)$ & $\rho_0 - A_0(T-T_m)$ & $T_m$--$T_b$ per metal & 0.3--3\% (Na), 0.7--0.8\% (Pb, LBE) & Sobolev (2020) Eq.\ [4], Table 7 \\
$\sigma(T)$ & $(\sigma_0 - A_0(T-T_m))\times10^{-3}$ & $T_m$--$T_b$ & 3--6\% (Na); Pb/LBE bands not independently confirmable from Sobolev's single collective figure & Eq.\ [11], Table 9 \\
$c_p(T)$ & $a+bT+cT^2+dT^{-2}$, molar, $/M$ & metal-specific window & 0--1\% (Na), 5--7\% (Pb, LBE) & Eq.\ [12], Table 10 \\
$h(T)$ & analytic integral of $c_p$ from $T_m$ & same as $c_p$ & inherited from $c_p$ & Eq.\ [14] --- \textbf{see the sign defect below} \\
$\mu(T)$ & $\mu_0\exp(E_0/RT)$ & metal-specific window & 5\% (Na, Pb), 7--10\% (LBE) & Eq.\ [17], Table 11 \\
$k(T)$ & metal-specific (Na, Pb linear in $T$; LBE quadratic) & metal-specific window & 0--8\% (Na, disputed), 0--15\% (Pb), 10--15\% (LBE) & Eq.\ [22], Table 13 \\
\end{longtable}

\textbf{Reference:} Sobolev, V. (2020), \emph{Properties of Liquid Metal Coolants: Na,
Pb, Pb-Bi}, Reference Module in Materials Science and Materials Engineering. Every
coefficient above was checked against this source and matches exactly, with two
exceptions, both reported rather than silently fixed (\code{CLAUDE.md}: ``change no
physics''):

\begin{itemize}
  \item \textbf{$h(T)$'s last term has the wrong sign, in all three metals.} Sobolev's
    own Eq.\ [14], the analytic integral of his Eq.\ [12] $c_p$, gives
    $d(1/T_m - 1/T)$ for the $d T^{-2}$ term; this module computes $d(1/T - 1/T_m)$
    instead, matching the IAEA handbook's form rather than Sobolev's. Confirmed by hand
    integration of the module's own $c_p$. Magnitude: roughly 1.1\% of Sodium's
    enthalpy rise at $T=T_b=1155$~K, growing with $|T-T_m|$. See
    \code{docs/OPEN\_QUESTIONS.md} Q35.
  \item \textbf{Three uncertainty bands do not match Sobolev's text.} \code{Sodium.uncert\_k
    = [0, 8\%]} where Sobolev states an examined spread of ``up to $\pm15\%$'';
    \code{Lead.uncert\_sig = [0, 5\%]} and \code{LBE.uncert\_sig = [0, 0.3\%]} where
    Sobolev gives only one collective ``(3--6)\%'' figure for all three coolants
    together. Not adjusted --- may trace to an unpublished-here source. See
    \code{docs/OPEN\_QUESTIONS.md} Q36.
\end{itemize}

Per \code{docs/DECISIONS.md} (``Scope: liquid metals''), only one correlation per
property is implemented; the CV's ``multiple literature-published correlations per
property'' claim is not supported by this repository.

\subsection{UO2}

\code{properties/matmod.py::UO2}. Fuel conductivity, conductivity integrals,
emissivity, thermal expansion, densification and swelling.

\subsubsection{Thermal conductivity}

Two models, both implemented.

\textbf{Klimenko-Zorin} (\code{k\_Klimenko}), the fresh-fuel model used through the
annular and rod SCA solvers:
\[
  \tau = T/1000, \qquad
  k(T) = \frac{100}{7.5408 + 17.692\tau + 3.6142\tau^2} + 6400\,\tau^{-5/2}\exp\!\left(-\frac{16.35}{\tau}\right)
\]
Its analytic antiderivative $\Theta_{Klimenko}(T) = \int k\,dT$ (\code{Theta\_Klimenko})
is also implemented, re-derived and numerically checked against
\code{scipy.integrate.quad} of \code{k\_Klimenko} itself (this is the function whose
erf-coefficient defect, D4 below, was fixed in Phase 2). \textbf{Valid range,
uncertainty and reference: \notest{}} --- checked against both
\code{docs/reference/MatLib\_Info.pdf} (PNNL-35702) and
\code{docs/reference/Hughes\_SCWR\_1.pdf} in full; neither names ``Klimenko-Zorin'' or
contains these coefficients, though Hughes (2014) Eq.\ (14) uses the same functional
family without naming its own source either.

\textbf{MATPRO / NFI, with burnup and gadolinia dependence} (\code{k\_NFI}):
\[
  f_{Bu} = 0.00187\,Bu, \quad g_{Bu} = 0.038\,Bu^{0.28}, \quad
  h_T = \frac{1}{1+396\exp(-Q/T)}
\]
\[
  k_{95} = \frac{1}{A + \alpha\,f_{gad} + BT + f_{Bu} + (1-0.9e^{-0.04Bu})\,g_{Bu}\,h_T} + \frac{E}{T^2}\exp\!\left(-\frac{F}{T}\right)
\]
with $A=0.0452$, $B=2.46\times10^{-4}$, $E=3.5\times10^9$, $F=16361$, $\alpha=1.1599$
(units as in the PDF). \textbf{Valid range:} $T\in[300,2800]$~K; gadolinia 0--10 wt\%;
burnup 0--90~GWd/MTU (UO2), 0--50~GWd/MTU (UO2-Gd$_2$O$_3$); density 90--98.6\%TD (not
itself an argument, so not checkable). \textbf{Uncertainty:} $\sigma=8.3\%$ (UO2),
$8.8\%$ (UO2-Gd$_2$O$_3$). \textbf{Reference:} Ohira \& Itagaki (1997, NFI model),
modified for burnup/gadolinia by Lanning et al.\ (2005), via PNNL-35702 section 2.1.1.1,
Eqs.\ 2-1/2-2. A fixed-step trapezoidal integral \code{Theta\_NFI} is also implemented
(no closed-form antiderivative exists for this model).

\textit{PDF discrepancy (Q17 item 1, resolved):} the PDF writes Klimenko's second term
as $6400/\tau^2\cdot\exp(-16.35/\tau)$; the code (and the published Fink/Klimenko-Zorin
form) has $6400\,\tau^{-5/2}$. The code is correct; this is a PDF transcription slip.

\subsubsection{Emissivity}

\code{eps(T)}: $\varepsilon = 0.7856 + 1.5263\times10^{-5}\,T$. \textbf{Valid range:}
gadolinia 0--10 wt\%, $T\in[300,2500]$~K. \textbf{Uncertainty:} $\sigma=0.072$,
absolute (PNNL-35702 section 2.1.5.3) --- this replaces an unsourced ``$\pm6.8\%$''
this docstring stated before Phase 3. \textbf{Reference:} PNNL-35702 section 2.1.5.1,
Eq.\ 2-12 (constant $0.78557$ there vs.\ $0.7856$ in the code; matches to the 4th
decimal, source of the small difference not established, preserved as found).

\textit{PDF discrepancy (Q17 item 2, resolved):} the PDF writes
$\varepsilon = 0.78557 + 1.5263\times T$, missing the \code{e-5}; without it emissivity
would exceed 1 at 300~K. The code is correct.

\subsubsection{Thermal Expansion}

\code{thrm\_expan(T)}: $\text{strain} = K_1 T - K_2 + K_3\exp(-E_d/(kT))$.
\textbf{Valid range:} UO2, UO2-Gd$_2$O$_3$, MOX; gadolinia 0--10 wt\%; $T\ge300$~K (no
enforced upper bound --- comparison data covers up to $\sim$3000~K).
\textbf{Uncertainty:} $\sigma=10.3\%$. \textbf{Reference:} PNNL-35702 section 2.1.4.1,
Eq.\ 2-9, Table 2-2.

\textit{PDF discrepancy (Q17 item 3, resolved):} the PDF's Eq.\ 4 has
$\exp(+1.32\times10^{-19}/kT)$ (positive exponent, which diverges as $T\to0$); the code
has $\exp(-E_d/kT)$, the correct MATPRO form.

\subsubsection{Swelling}

\notest{} in the PDF's own text (the PDF's TOC entry has no body); the code is
\emph{ahead} of the manual here. Four functions implemented:
\code{swelling\_solid(Bu, gad)} (piecewise linear in burnup, gadolinia branch
$0.0005\,Bu$; PNNL-35702 section 2.1.8.1, Eqs.\ 2-17/2-18 --- note PNNL-35702 itself
gives two contradictory uncertainty figures for the same stated burnup condition,
Q37), \code{swelling\_gas(T, Bu)} (active only for $Bu\ge40$~GWd/MTU and
$1233\text{--}2105$~K, section 2.1.8.1 Eqs.\ 2-19--2-21), \code{densification(T, Bu,
rho\_TD, T\_sint)} and its two helpers \code{dens\_max}/\code{dens\_B} (Rolstad et al.\
1974, via PNNL-35702 section 2.1.7.1, Eqs.\ 2-15/2-16 --- PNNL-35702 itself states
uncertainty ``could not be meaningfully established'' for this model due to data
scatter).

\subsection{Steels}

\subsubsection{HT9 Steel Models}

\code{properties/matmod.py::HT9}. Substantial and, before this phase, undocumented in
the manual (\notest{} in the PDF's own text, same as UO2 swelling above). Thermal
conductivity ($k=A_0+A_1T$, 293--873~K, Akiyama 1991), specific heat (piecewise linear
with a break at 800.15~K, 298--873~K, Yamanouchi 1992), melting temperature (constant
973~K, the HT-9/metallic-fuel eutectic, Baker \& Wilson 1992), density (constant
7750~kg/m$^3$, Akiyama 1991), emissivity (constant 0.9, Dutt \& Baker 1974), thermal
expansion ($A_1+A_2T+A_3T^2$, 298.15--1073.15~K, Yamanouchi 1992 --- see the units note
below), Young's and shear moduli ($E,G = A_0+A_1T$, with \textbf{$T$ in degrees Celsius,
not kelvin, unlike every other temperature argument in this module} --- confirmed
against PNNL-35702's own ``Where,'' clause, so this is not a transcription artifact),
Meyer's hardness (delegates to the zirconium-based correlation below 1235~K, per
PNNL-35702 section 3.3.9's own instruction), primary/secondary/tertiary thermal creep
strain rate and irradiation creep strain rate (Akiyama 1991, section 3.3.10, all
298.15--873.15~K, no stated uncertainty for any of the strain-rate equations), and
yield stress (cubic in $T$, Akiyama 1991, section 3.3.11).

\textit{Units flag (not resolved, not a code change):} PNNL-35702 labels
\code{thrm\_expan}'s output ``$\alpha$ = thermal expansion coefficient, K$^{-1}$'', but
this module's docstring calls it a dimensionless ``strain'', matching how the
Zircaloy and UO2 expansion functions are used elsewhere. Whether the fitted
coefficients are actually a CTE or a strain under a borrowed symbol is not resolvable
from the formula alone; see \code{docs/OPEN\_QUESTIONS.md} Q40.

\textbf{Reference (all of HT9, unless noted otherwise above):} PNNL-35702 (Geelhood et
al.\ 2024), \emph{MatLib-1.2.1: Nuclear Material Properties Library}, section 3.3.

\subsubsection*{D9 stainless steel (not in the PDF TOC, but a named section 1.4 stub)}

\code{properties/matmod.py::D9\_SS}. Two constants only, deliberately not a
temperature-dependent model since no source supplies one: $k=18.9$~W/m-K (evaluated
once at 650~K, Hughes' average clad temperature, held constant; \code{T} is accepted
and ignored) and $\rho=8100$~kg/m$^3$. \textbf{Reference:} Leibowitz \& Blomquist
(1988), via Hughes, Pelaez, Schubring \& Jordan, \emph{Nucl.\ Eng.\ Des.} 270 (2014)
412--420, section 4 and Table 1.

\subsubsection*{Zircaloy (not in the PDF TOC at all)}

\code{properties/matmod.py::Zircalloy}. The largest single block of code in this
module: thermal conductivity ($k$, cubic below 2098~K then a 36.0~W/m-K phase-transition
plateau, PNNL-35702 section 3.1.1.1 Eqs.\ 3-1/3-2), specific heat (14-point
piecewise-linear table, section 3.1.2.1 Table 3-2), melting temperature (constant
2123.15~K), density (constant 6520~kg/m$^3$), axial and diametral thermal expansion
(piecewise, with a 21-point table across the alpha-to-beta transition, section 3.1.4.1
Eqs.\ 3-4/3-5), emissivity (oxide-thickness- and temperature-dependent, section 3.1.5.1
Eqs.\ 3-6/3-7), Young's and shear moduli (temperature/cold-work/fluence/oxide-dependent,
section 3.1.7.1 Eqs.\ 3-9/3-10), Meyer's hardness (cubic-in-$T$ below 1235~K then a
plateau, section 3.1.8.1 Eq.\ 3-11), axial irradiation growth (power-law in fluence,
alloy-specific coefficients, section 3.1.9.1 Eqs.\ 3-12--3-15), effective (von Mises)
stress from thick-wall Lam\'e principal stresses (section 3.1.10.1 Eqs.\ 3-23--3-26),
and a full thermal + irradiation creep-strain-rate model (section 3.1.10.1 Eqs.\
3-16--3-19) covering Zircaloy-2, Zircaloy-4, M5, ZIRLO and Optimized ZIRLO. All from
PNNL-35702, filled in during Phase 3.

\textit{Documentation quirks found while filling this section, not code defects (see
\code{docs/OPEN\_QUESTIONS.md} for detail):} PNNL-35702 states two different
correlation-family assignments for Optimized ZIRLO's creep in two places in the same
section (Q38); its irradiation creep-rate flux range is stated four orders of magnitude
off from its own formula's stated unit (Q39); and its solid-swelling uncertainty
section states two different sigma values for the same burnup condition (Q37, this
applies to UO2, listed above, not Zircaloy).

\subsection{Gas}

\subsubsection{Thermal Conductivity}

\code{properties/matmod.py::Gas.k(gas, T)}: $k = A\,T^B$, tabulated per species (He, Ar,
Kr, Xe, H$_2$, N$_2$, Air). \textbf{Valid range:} He/Ar/N$_2$ 273--2500~K, Kr
273--2300~K, Xe 273--2200~K, H$_2$ 273--2000~K; \textbf{Air's range and uncertainty are
\notest{}} --- present in PNNL-35702's fitting-constant table (section 4.1.1) but absent
from its applicability/uncertainty bullets (section 4.1.3). \textbf{Uncertainty
(absolute):} He $8.99\times10^{-3}$, Ar $9.66\times10^{-4}$, Kr $8.86\times10^{-4}$, Xe
$5.34\times10^{-4}$, H$_2$ $1.67\times10^{-2}$, N$_2$ $1.99\times10^{-3}$~W/m-K.
\textbf{Reference:} PNNL-35702 section 4.1.1, Eq.\ 4-1, Table 4-1. This feeds
\code{pin/gap.py::htc\_gap}'s fill-gas conductivity; the gap conductance model's
composite conduction+radiation form itself traces to Von Ubisch et al.\ (1958), per
\code{docs/PHYSICS\_REVIEW.md}.

\subsection{Future materials}

\notimpl{} A CO$_2$ equation of state and HTGR gaseous-coolant property correlations
were explicitly future work in the original PDF, and remain unimplemented. Out of
scope per \code{docs/GAP\_ANALYSIS.md} Part D.

% =====================================================================================
\section{Heat Transfer Coefficient Models}
% =====================================================================================

\code{correlations/htc.py}. Flat-namespace classes \code{Water}, \code{SCW},
\code{Sodium}, \code{Lead}. Every correlation that needs a wall temperature solved from
a heat flux is wrapped by a torchsolve-backed bracketed root-find
(\code{\_solve\_Tw}/\code{\_solve\_Tw\_scw}) --- \code{scipy.optimize} is retired from
this library by decision (\code{docs/DECISIONS.md}, ``Root finding'').

\subsection{Single Phase Water}

\subsubsection{Dittus Boelter}

\code{Water.Dittus(Props, G, D)}:
\[
  Pr = \frac{\mu c_p}{k}, \quad Re = \frac{GD}{\mu}, \quad Nu = 0.023\,Re^{0.8}Pr^{0.4},
  \quad htc = \frac{Nu\,k}{D}
\]
This implementation always uses the heating exponent $n=0.4$ (every call site in this
repository heats the coolant); a cooling variant with $n=0.3$ would be a new capability,
not a preserved one, and is not added. \textbf{Valid range:} $Re>10{,}000$;
$0.7<Pr<160$; $L/D>10$; small bulk-to-wall $\Delta T$. \textbf{Uncertainty:} roughly
$\pm25$--$45\%$. \textbf{Reference:} Dittus, F.W. and Boelter, L.M.K., \emph{Univ.\
California Publ.\ Eng.}, 2:443 (1930).

\textbf{Physics fix (D1, commit \code{cd9c376}):} before this cleanup, this function
used a leading constant of \textbf{0.026} (Colburn's constant, which belongs to a
different, $Pr^{1/3}$-exponent correlation), not the published \textbf{0.023}. At
$Re=1.11\times10^5$, $Pr=0.825$, that returned 15{,}715~W/m$^2$-K against a correct
13{,}902~W/m$^2$-K --- \textbf{13.0\% high}, propagating into every single-phase result
that used it. Corrected to 0.023.

\subsubsection{Petchukov}

\code{Water.Petchukov(Props, G, D)} (Petukhov, spelling preserved from the original
source throughout):
\[
  f = (0.790\ln Re - 1.64)^{-2}, \qquad
  Nu = \frac{(f/8)\,Re\,Pr}{1.07 + 12.7\sqrt{f/8}\,(Pr^{2/3}-1)}, \qquad htc=\frac{Nu\,k}{D}
\]
\textbf{Valid range:} $10^4<Re<5\times10^6$; $0.5<Pr<2000$. \textbf{Uncertainty:}
$\pm5\%$. \textbf{Reference:} Petukhov, B.S.\ (1970).

\subsubsection{Gnielinski}

\code{Water.Gnielinski(Props, G, D)}:
\[
  f = (1.82\log_{10}Re - 1.64)^{-2}, \qquad
  Nu = \frac{(f/8)(Re-1000)Pr}{1+12.7\sqrt{f/8}(Pr^{2/3}-1)}, \qquad htc=\frac{Nu\,k}{D}
\]
\textbf{Valid range:} $2300<Re<5\times10^6$; $0.5<Pr<2000$ --- valid to a lower
Reynolds number than Petukhov, useful near the laminar-turbulent transition.
\textbf{Uncertainty:} $\pm10\%$. \textbf{Reference:} Gnielinski, V.\ (1976).

\subsection{Two Phase Water}

\subsubsection{Schrock and Grossman}

\code{Water.SchrockGrossman(Props\_l, Props\_v, htc\_lo, x, G, D, q\_pp)}:
\[
  X_{tt} = \left(\frac{\mu_l}{\mu_v}\right)^{0.1}\left(\frac{\rho_v}{\rho_l}\right)^{0.5}\left(\frac{1-x}{x}\right)^{0.9},
  \qquad h_{tp} = htc_{lo}\left(1.11\,X_{tt}^{-0.66} + \frac{7400\,q''}{G\,h_{fg}}\right)
\]
\textbf{Reference:} Schrock, V.E.\ and Grossman, L.M.\ (1959). \textbf{Valid range /
uncertainty:} \notest{} range; uncertainty $\pm30\%$.

\textit{Documentation fix (D12), not a physics change:} this function's docstring, before
Phase 2, described a different formula ($h_{tp}=2.5\,h_l(1/X_{tt})^{0.75}$) than the
body actually computed --- an apparent transcription error copied from a different
correlation. The body's formula (above) is unchanged and matches
\code{docs/reference\_code/SCA\_Example.py}'s \code{htc2phi} and the standard published
Schrock-Grossman form; only the docstring was corrected to describe what the code
already did.

\subsubsection{Chen}

\code{Water.Chen\_H2O\_dT} (explicit, given $T_w$) and \code{Water.Chen\_H2O} (implicit,
solves for $T_w$ given a heat flux):
\[
  h_{tp} = F h_c + S h_{nb}, \qquad h_c = 0.023\,Re_l^{0.8}Pr_l^{0.4}\frac{k_l}{D}
\]
\[
  F = \begin{cases}1.0 & 1/X_{tt}\le0.1\\ 2.35\,(1/X_{tt}+0.213)^{0.736} & \text{otherwise}\end{cases}
  \qquad (\text{a batched \code{xp.where}, not a Python \code{if}})
\]
\[
  h_{nb} = 0.00122\,\frac{k_l^{0.79}c_{p,l}^{0.45}\rho_l^{0.49}}{\sigma^{0.5}\mu_l^{0.29}h_{fg}^{0.24}\rho_v^{0.24}}\,\Delta T_{sat}^{0.24}\,\Delta P_{sat}^{0.75},
  \qquad S = \frac{1}{1+2.53\times10^{-6}Re_{tp}^{1.17}},\ \ Re_{tp}=Re_l F^{1.25}
\]
\textbf{Reference:} Chen, J.C.\ (1966). \textbf{Valid range / uncertainty:} \notest{}
range; uncertainty $\pm30\%$.

\subsubsection{Bjorge}

\code{Water.Bjorge\_dT} / \code{Water.Bjorge}: the same convective/nucleate-boiling
ingredients as Chen, combined as a root-sum-square (asymptotic) blend instead of Chen's
suppression-factor superposition:
\[
  h_{tp} = \sqrt{h_{fc}^2 + h_{nb}^2}, \qquad h_{fc} = F\,h_l
\]
with $F$ and $h_{nb}$ the same functional forms as Chen's above. \textbf{Reference:}
Bjorge, R.W., Hall, G.R.\ and Rohsenow, W.M.\ (1982). \textbf{Uncertainty:} $\pm25\%$.

\subsection{Supercritical Water}

\subsubsection{Swenson Correlation}

\code{SCW.Swenson\_dT} (explicit) / \code{SCW.Swenson} (implicit, non-monotone-aware
solve --- \code{htc}$(T_w)$ peaks near the pseudocritical temperature, which
\code{torchsolve}'s own README names this correlation by name as its motivating case):
\[
  \overline{c_p} = \frac{h_w-h_b}{T_w-T_b}, \qquad A = \frac{\overline{c_p}}{c_{p,w}},
  \qquad B = \frac{\rho_w}{\rho_b}
\]
\[
  Nu_w = 0.00459\,Re_w^{0.923}\,Pr_w^{0.613}\,A^{0.613}\,B^{0.231}, \qquad htc = \frac{Nu_w k_w}{D}
\]
\textbf{Valid range / uncertainty:} \notest{} range; $\pm25\%$. \textbf{Reference:}
Swenson, H.S., Carver, J.R.\ and Kakarala, C.R.\ (1965), per Hughes, Pelaez, Schubring
\& Jordan, \emph{Nucl.\ Eng.\ Des.} 270 (2014) 412--420, Eq.\ (8).

\textbf{Physics fix (D2, commit \code{7adc1ac}):} before this cleanup, the averaged
heat-capacity ratio was computed as $\overline{c_p}/c_{p,b}$ (bulk-referenced) with
exponent \textbf{0.231} on that factor; the published Swenson $Pr_{\overline{bar},w}$
grouping requires the ratio referenced to the \emph{wall} ($\overline{c_p}/c_{p,w}$,
as above) with exponent \textbf{0.613}. Bulk-referencing leaves a spurious
$c_{p,w}/c_{p,b}$ factor in the result. \textbf{14.4\% high at a representative state.}
Confirmed against Hughes et al.\ (2014) Eq.\ (8) once that source was supplied
(\code{docs/reference/Hughes\_SCWR\_1.pdf}); the 0.231 that had been on the wrong factor
is, not coincidentally, the lead constant of the neighboring Bishop correlation on the
same page of that source (Q29) --- likely the origin of the mix-up. Corrected in both
live copies (\code{correlations/htc.py} and \code{sca/rod.py::Swenson}).

\subsubsection{Chen et al}

\code{SCW.Chen\_SCW\_dT} / \code{SCW.Chen\_SCW}. Per \code{docs/DECISIONS.md}, promoted
to a first-class supercritical option and the \textbf{recommended default}
(Swenson above is retained as the legacy benchmark):
\[
  Nu_b = 0.46\,Re_b^{0.16}\left(\frac{Pr_w}{Pr_b}\right)^{0.1}\left(\frac{\nu_w}{\nu_b}\right)^{-0.55}\left(\frac{\overline{c_p}}{c_{p,b}}\right)^{0.88}\left(\frac{Gr_b^*}{Gr_b}\right)^{0.81}
\]
where the heat-flux-to-buoyancy Grashof ratio collapses, under the Boussinesq
approximation $\beta_b=(\rho_b-\rho_w)/(\rho_b(T_w-T_b))$, to $qD/(k_b(T_w-T_b))$ ---
$g$, $\beta$ and $\nu_b$ all cancel algebraically. \textbf{Valid range:} bulk enthalpy
278--3169~kJ/kg; $G$ 201--2500~kg/m$^2$-s; $q$ 129--1735~kW/m$^2$; $P$
22--34.3~MPa (not itself an argument, so not checked); $D$ 6--26~mm.
\textbf{Uncertainty:} MAD 5.4\%; 95.7\% of the 5366-point regression database within
$\pm15\%$. \textbf{Reference:} Chen, W.\ and Fang, X.\ (2014), \emph{Int.\ J.\ Heat Mass
Transfer} 78, 156--160 --- the source PDF is in \code{docs/reference/}, and this
function's docstring is the most thoroughly documented in the codebase (range, MAD,
database size, and the Boussinesq derivation, all stated), used as the docstring
exemplar for the rest of the cleanup.

\subsection{Liquid Sodium}

\subsubsection{Notter-Sleicher correlation}

\notimpl{} No source for the Notter-Sleicher correlation exists anywhere in this
repository --- checked against \code{docs/reference/MatLib\_Info.pdf}, Todreas \&
Kazimi Volume 1, the Cheng SCWR review, the Meyer heat-transfer-coefficient review, and
the Wu paper; none mention it. Per the owner's own rule (implement only what can be
implemented accurately), it is left as a documented gap rather than guessed. See
\code{docs/OPEN\_QUESTIONS.md} Q15.

\subsubsection*{What is implemented instead (not in the PDF TOC)}

Two classic sodium correlations for the pair of standard boundary conditions, both from
Todreas \& Kazimi \emph{Nuclear Systems Volume 1}, 3rd ed., section 10.5.3.1:

\code{Sodium.Lyon} (constant heat flux --- the boundary condition a fuel pin actually
approximates):
\[
  Pe = \frac{GDc_p}{k}, \qquad Nu = 7.0 + 0.025\,Pe^{0.8}, \qquad htc=\frac{Nu\,k}{D}
\]
\code{Sodium.SebanShimazaki} (uniform wall temperature): the same form with the
conduction-floor constant 5.0 in place of 7.0. Neither correlation has a stated Peclet
range or uncertainty band in T\&K.

\code{Sodium.Mikityuk} --- the rod-\emph{bundle} liquid-metal correlation (distinct from
the two circular-tube correlations above; it returns a Nusselt number outright rather
than a tube-to-bundle correction factor, which is why it lives here and not in
\code{correlations/bundle.py}):
\[
  x = \frac{pitch}{D}, \qquad Nu = 0.047\left(1-e^{-3.8(x-1)}\right)\left(Pe^{0.77}+250\right), \qquad htc=\frac{Nu\,k}{D}
\]
\textbf{Valid range:} $30\le Pe\le5000$; $1.1\le pitch/D\le1.95$. \textbf{Uncertainty:}
MAE $-0.1$, RMS $1.9$ (absolute Nusselt-number units, not relative). \textbf{Reference:}
Mikityuk, K.\ (2009), \emph{Nucl.\ Eng.\ Des.} 239, 680--687, Eq.\ (14). T\&K's own
Eq.\ (10.133) mis-typesets the Peclet exponent as ``$(Pe<0.77+250)$'' --- confirmed by
rendering the source page as an image; the ``$<$'' is a mis-typeset superscript. Resolved
against Mikityuk's own paper once supplied; see \code{docs/OPEN\_QUESTIONS.md} Q41.

\subsection{Liquid Lead and Lead-Bismuth Eutectic}

\subsubsection{Shen}

\code{Lead.Shen(Props, T, G, D)}:
\[
  Pe = Re\,Pr, \qquad Nu = 10.287\,Pe^{+0.1175} + \frac{0.0599}{2.5}\,Pe^{0.7575}, \qquad htc=\frac{Nu\,k}{D}
\]
\textbf{Valid range / uncertainty / reference: \notest{}.} \textbf{Unresolved sign
(D3):} the pre-cleanup codebase had three copies of this correlation, and two of the
three used a \emph{negative} exponent on the leading Peclet term
($Nu=10.287\,Pe^{-0.1175}+\ldots$) --- at $Pe=500$ that is $Nu=24.0$ vs.\ $Nu=7.6$, a
substantial difference. No source document in this repository resolves which sign is
correct. This module keeps the positive-exponent form it already had; per
\code{docs/DECISIONS.md} (``Scope: lead is a toy''), the lead channel in the annular SCA
files is an experiment, not an active path, so this is property-library infrastructure
and the Phase 7 uncertainty figure only. See \code{docs/OPEN\_QUESTIONS.md} Q13.

% =====================================================================================
\section{Friction Models}
% =====================================================================================

\code{correlations/friction.py}. Flat-namespace classes \code{f\_water}, \code{f\_SCW},
\code{Spacer}.

\subsection{Single Phase Water}

\subsubsection{McAdams}

\code{f\_water.McAdams(Props, G, D)}: $Re=GD/\mu$, $f=0.184\,Re^{-0.2}$.
\textbf{Valid range:} $30{,}000<Re<1{,}000{,}000$ (as stated in the original source).
\textbf{Uncertainty / reference: \notest{}.}

\subsubsection*{Blasius (not in the PDF TOC)}

\code{f\_water.Blasius(Props, G, D)}: $f=0.316\,Re^{-0.25}$. \textbf{Valid range:}
$Re<10^5$. The lower-Reynolds-number companion to McAdams. \textbf{Uncertainty /
reference: \notest{}.}

\subsubsection{Colebrook}

\code{f\_water.Colebrook(Props, G, D, roughness, n\_iter, tol, return\_convergence,
check\_range)} --- \notimpl{} before Phase 4, now implemented:
\[
  \frac{1}{\sqrt{f}} = -2\log_{10}\left(\frac{roughness/D}{3.70} + \frac{2.51}{Re\sqrt{f}}\right)
\]
Solved by a fixed-iteration-count \emph{fixed-point} substitution $x=1/\sqrt f$ (not
Newton): the map is strongly contracting for turbulent $x$ (contraction factor roughly
0.07--0.17), branch-free, and keeps every iterate positive because of the logarithm,
whereas Newton could leave the physical branch on a bad initial guess. Seeded from the
Haaland approximation (within 2\% of Colebrook per T\&K), so the loop starts close.
Vectorized over the whole batch with a fixed iteration count, differentiable end to
end. \textbf{Valid range:} turbulent flow, $Re\gtrsim4000$. \textbf{Uncertainty:}
Colebrook is itself the reference the explicit fits (McAdams, Blasius) are measured
against, so no percentage uncertainty is stated. \textbf{Reference:} Colebrook, C.F.\
(1939), \emph{J.\ Inst.\ Civ.\ Eng.}, as given by Todreas \& Kazimi \emph{Nuclear
Systems Volume 1}, 3rd ed., Eq.\ (9.89), in this module's Darcy convention throughout.

\subsection{Supercritical Water}

\subsubsection{Filoneko Correlation}

\code{f\_SCW.Filonenko(Props, G, D, Props\_w=None)} (Filonenko, spelling as in the PDF):
\[
  f = \frac{1}{(1.82\log_{10}Re - 1.64)^2}
\]
Algebraically the isothermal part of Hughes et al.\ (2014) Eq.\ (9) (Petrov \& Popov
1988): $1.82\log_{10}Re - 1.6437$ there vs.\ $-1.64$ here. An optional Petrov-Popov
supercritical density correction is available, \textbf{defaulting off} so existing call
sites keep the isothermal behavior they were validated with (\code{docs/DECISIONS.md}):
supplying \code{Props\_w} (wall-state properties) multiplies in
$(\rho_w/\rho_b)^{0.4}$. This runs on both channels of the annular SCA per
\code{docs/DECISIONS.md} (``Wu friction''). \textbf{Valid range / uncertainty:}
\notest{}. \textbf{Reference:} Petrov \& Popov (1988), as cited by Hughes et al.\ (2014)
Eq.\ (9), for the isothermal form implemented here.

\subsubsection{Wu Correlation for Rod Bundles}

\code{f\_SCW.Wu(Props, G, D)}:
\[
  f = 0.014\,f_{Filonenko}^{-0.12}\,Pr^{-0.23}
\]
\textbf{Not used in the annular SCA path today} (Filonenko runs on both channels per
owner instruction); range-limited to $G\le1000$~kg/m$^2$-s in \code{ranges.py} pending
the owner's own high-mass-flux correction factor. \textbf{Uncertainty / reference:
\notest{}} --- the negative exponent on the isothermal friction factor is unusual for a
bundle correction and is unconfirmed against any source. See
\code{docs/OPEN\_QUESTIONS.md} Q14.

\subsubsection*{Spacer-grid loss (not in the PDF TOC)}

\code{Spacer.blah2()}: \notimpl{}, and honestly so --- converted from a bare
\code{return} (which silently did nothing) to \code{raise NotImplementedError}. No
formula, source, or even a description of what this stub was meant to compute survives
anywhere in the repository beyond the class/method name suggesting a spacer-grid
friction or mixing correction. See \code{docs/OPEN\_QUESTIONS.md} Q33.

% =====================================================================================
\section{Fuel Pin Heat Transfer Models/Routines}
% =====================================================================================

\code{pinthac/pin/}: \code{gap.py}, \code{clad.py}, \code{cylindrical.py},
\code{annular.py}.

\subsection{Gap Heat Transfer}

\code{pin/gap.py::htc\_gap(Tfo, Tci, delta, kgas, eps\_c=1.0, eps\_f=1.0, units='J')}:
conduction plus radiation across an open (not mechanically closed) gas gap, treated as a
plane layer (valid when the diametral gap is small compared to the fuel radius):
\[
  T_{ave} = \tfrac12(T_{fo}+T_{ci}), \qquad h_{cond} = \frac{k_{gas}(T_{ave})}{\delta}
\]
\[
  h_{rad} = \frac{\sigma}{\varepsilon_f^{-1}+\varepsilon_c^{-1}-1}\,(T_{fo}^2+T_{ci}^2)(T_{fo}+T_{ci})
\]
(algebraically $\sigma(T_{fo}^4-T_{ci}^4)/(\varepsilon^{-1}\text{-sum}\cdot(T_{fo}-T_{ci}))$,
factored to stay finite as $T_{fo}\to T_{ci}$), $h_{gap}=h_{cond}+h_{rad}$. Does NOT
include mechanical contact-closure conductance. \textbf{Valid range / uncertainty /
reference for the composite form: \notest{}}; the gas conductivity usually passed in
(\code{MatMod.Gas.k}) is Von Ubisch et al.\ (1958), per \code{docs/PHYSICS\_REVIEW.md}.

\textbf{Physics fix (D5, commit \code{7f42d80}):} a duplicate hand-rolled copy of this
formula inside \code{SCW\_Pb\_Ann\_SCA.py} used gas conductivity
$k_{gas}=15.8\times10^{-4}\,T^{-0.79}$ --- a \textbf{negative} exponent, where the other
two copies (including this one) used $+0.79$. At 600~K that puts $k_{gas}$ at
$\sim10^{-5}$~W/m-K instead of $\sim0.25$~W/m-K, four orders of magnitude low, which
effectively zeroes conduction across the gap and leaves radiation to carry the whole
transfer alone. Corrected to $+0.79$; \code{pin/gap.py::htc\_gap} (this function, which
already had the correct sign and additionally carries the emissivity factor the
duplicate lacked) is adopted as the one canonical implementation.

\subsection{Clad conduction}

\code{pin/clad.py::T\_ci(Rco, Rci, kc, Tco, qp)}:
\[
  T_{ci} = T_{co} + \frac{qp}{2\pi k_c}\ln\!\left(\frac{R_{co}}{R_{ci}}\right)
\]
Handles the annulus case, where the \emph{inner} channel's clad has $R_{co}<R_{ci}$ by
naming convention: the larger of the two radii goes on top of the log ratio either way,
by flipping the sign rather than the ratio. Assumes constant $k_c$ across the wall step
(a caller wanting temperature-dependent conductivity evaluates $k_c$ at the layer's mean
temperature first, as \code{sca/annular.py::closure} does). \textbf{Valid range /
uncertainty / reference: \notest{}} --- standard steady-state cylindrical (Fourier)
conduction; no fitted database to be out of range of.

\subsection{Cylindrical Heat transfer}

\code{pin/cylindrical.py}. The manual's description --- ``a numerical solver ... solved
iteratively for $T$ given the radius'' --- is implemented by \code{Cyl\_T}, with three
supporting pieces:

\begin{itemize}
  \item \code{Cyl\_HT(r, q\_vol, kint, C)}: the \textbf{constant-conductivity} special
    case, $T(r) = C - q_{vol}r^2/(4\,kint)$, sign confirmed against
    \code{Annular\_Heat\_Transfer\_Final.pdf} Eq.\ (2) with $r_i\to0$ (see also the PDF
    discrepancy below).
  \item \code{Cyl\_Theta(r, rfo, q3, Theta\_fo)}: the conductivity-integral value at
    radius $r$, given its value at the pellet surface --- exact and non-iterative for a
    solid pellet, since the general annular two-constant solution's $\ln(r)$ term
    diverges at the centerline and forces that constant to zero: $\Theta(r) = \Theta_{fo}
    + \tfrac{q_3}{4}(r_{fo}^2-r^2)$.
  \item \code{Cyl\_qpp(r, q3)} and \code{Cyl\_qlin(r, q3)}: the radial heat flux
    ($q_3 r/2$) and linear heat rate inside radius $r$ ($\pi q_3 r^2$), from geometry
    alone, needing no conductivity model.
  \item \code{Cyl\_T(r, rfo, q3, Theta\_fo, Theta\_func, k\_func, T\_lo, T\_hi)}: the
    actual iterative solve --- computes the target $\Theta$ via \code{Cyl\_Theta}, then
    inverts $\Theta_{func}(T)=\Theta_{target}$ for $T$ by bisection-then-Newton
    (\code{pinthac.solvers.bisect\_newton}, not \code{torchsolve}'s bracket-guarded
    solver, because $\Theta$ is strictly increasing in $T$ --- $d\Theta/dT=k(T)>0$ for
    any physical conductivity, so the root is unique and there is no non-monotone
    branch to guard against).
\end{itemize}

\textbf{Valid range / uncertainty:} not applicable to the geometric identities
(\code{Cyl\_qpp}/\code{Cyl\_qlin}, exact); inherited from the conductivity model for
\code{Cyl\_T}. \textbf{Reference:} this section of the manual itself, and
\code{Annular\_Heat\_Transfer\_Final.pdf} Eq.\ (2) with $r_i\to0$.

\textit{PDF discrepancy (Q17 item 4, resolved):} the PDF's section 4.3 writes the
cylindrical solve as $\int k\,dT = +q'''r^2/4+\ldots$; both the code and
\code{Annular\_Heat\_Transfer\_Final.pdf} Eq.\ 2 have $-q'''r^2/4$. The code is correct.

\subsection{Annular Heat Equation}

\code{pin/annular.py}. Implements the Kirchhoff-transformed radial solve of
\code{Annular\_Heat\_Transfer\_Final.pdf} section 1.1: substituting
$\Theta[k_f](T(r))=\int k_f\,dT$ for $T(r)$ linearizes the otherwise-nonlinear annular
conduction equation.

\textbf{\code{Ann\_Theta}} builds $\Theta[k_f]$ as a cumulative-trapezoid numerical
integral, returned as a SciPy \code{interp1d} --- \textbf{not differentiable}, a known
architectural limitation carried over from the pre-cleanup code (silently strips
gradients from a torch input passed through it, per \code{docs/AUDIT.md}), not fixed in
this cleanup.

\textbf{\code{Ann\_HT}(ri, ro, q3, Theta\_i, Theta\_o)} solves the closed-form
two-constant boundary value problem for $r_i\le r\le r_o$ with uniform volumetric
generation $q_3$:
\[
  \Theta[k_f](T(r)) = -\frac{q_3}{4}r^2 + C_1\ln r + C_2
\]
\[
  \alpha_i=\Theta_i+\tfrac{q_3}{4}r_i^2,\quad \alpha_o=\Theta_o+\tfrac{q_3}{4}r_o^2,
  \qquad C_1=\frac{\alpha_i-\alpha_o}{\ln r_i-\ln r_o},\quad
  C_2=\frac{(\ln r_i)\alpha_o-(\ln r_o)\alpha_i}{\ln r_i-\ln r_o}
\]
\textbf{\code{Ann\_qpp}(r, q3, C1)}: $q''(r)=\tfrac{q_3}{2}r - C_1/r$, in the signed
$+r$ Fourier convention.

\textbf{Sign convention (Q18, resolved by \code{docs/DECISIONS.md}):} the reference
PDF's $T_{fo}(r_j)=T_{m,j}+q''_j/htc_j$ is only correct at the \emph{outer} surface. The
inner coolant sits on the $-r$ side of the fuel, so the physically-outward flux there is
the \emph{negative} of \code{Ann\_qpp}'s signed function --- verified two ways (the
energy-balance invariant below, and both walls coming out hotter than their coolants,
which the PDF's literal convention does not give) plus an independent \code{sympy}
solve. Callers use $T_{fo}(r_i) = T_{m,i} - q''_i/htc_i$ (a minus) at the inner surface.
This is a deliberate, documented departure from the PDF's literal wording, approved by
the owner.

\subsubsection*{The iteration: \code{Ann\_flux\_split} (Figure 2 of the derivation PDF)}

\code{Ann\_flux\_split(ri, ro, q\_lin, Tm\_i, Tm\_o, htc\_i, htc\_o, Theta\_func,
f\_prev=None, n\_iter, tol)} closes the circular dependency \code{Ann\_HT} leaves open:
the fuel surface temperatures it needs depend on the surface fluxes, which are what the
conduction solve produces. Five repeated steps, matching the derivation PDF's Figure 2
exactly:

\begin{enumerate}
  \item From the current linear-heat split, surface fluxes:
    $q''_o = q'_o/(2\pi r_o)$, $\ q''_i = q'_i/(2\pi r_i)$.
  \item Fuel surface temperatures, across the coolant-to-surface resistance:
    $T_{fo}(r_o)=T_{m,o}+q''_o/htc_o$, $\ T_{fo}(r_i)=T_{m,i}-q''_i/htc_i$ (the sign
    convention above).
  \item The conductivity integral, evaluated \textbf{forward only}:
    $\Theta_j=\Theta[k_f](T_{fo}(r_j))$ --- never inverted, which is the entire point of
    the scheme: an inverse would need a Newton solve per surface per iteration per
    axial node, and $\Theta$ is only available as a forward integral of $k_f$.
  \item Cramer's rule on the resulting $2\times2$ linear system (this is \code{Ann\_HT}
    itself): $C_1 = (\alpha_i-\alpha_o)/(\ln r_i - \ln r_o)$.
  \item Updated fluxes from $C_1$ alone (\code{Ann\_qpp}; $C_2$ is not needed here).
\end{enumerate}

repeated until the surface fluxes stop moving. An optional warm start
(\code{f\_prev}, the previous axial node's heat ratio $q'_i/q'_o$) is the option
described in the PDF's section 1.2 --- since the split varies slowly along a channel,
starting from the previous node's answer typically needs only a small correction.

\begin{center}
\begin{tikzpicture}[node distance=1.05cm and 2.6cm, >=Stealth, font=\small]
  \node (start) [draw, rounded corners, fill=gray!10] {linear-heat split $q'_i, q'_o$};
  \node (step1) [draw, below=of start] {1. surface fluxes $q''_i, q''_o$};
  \node (step2) [draw, below=of step1] {2. fuel surface temps $T_{fo}(r_i), T_{fo}(r_o)$};
  \node (step3) [draw, below=of step2] {3. $\Theta_i, \Theta_o = \Theta[k_f](T_{fo})$ (forward only)};
  \node (step4) [draw, below=of step3] {4. Cramer's rule $\to C_1$ (\code{Ann\_HT})};
  \node (step5) [draw, below=of step4] {5. updated $q'_i, q'_o$ from $C_1$ (\code{Ann\_qpp})};
  \node (check) [draw, diamond, aspect=2.2, below=of step5, fill=gray!10] {converged?};
  \node (done) [draw, rounded corners, below=of check, fill=gray!10] {return $q_i, q_o, T_{fo,i}, T_{fo,o}$};
  \draw[->] (start) -- (step1);
  \draw[->] (step1) -- (step2);
  \draw[->] (step2) -- (step3);
  \draw[->] (step3) -- (step4);
  \draw[->] (step4) -- (step5);
  \draw[->] (step5) -- (check);
  \draw[->] (check) -- node[right]{yes} (done);
  \draw[->] (check.west) to[out=180,in=180] node[left]{no} (step1.west);
\end{tikzpicture}
\end{center}

\textbf{Energy balance is exact at every iterate, not just at convergence} ---
$q'_i+q'_o = \pi q_3 (r_o^2-r_i^2)$ with the $C_1$ terms cancelling identically. This
iteration never has to converge the total power, only how it splits, which is why an
unconverged result still conserves energy and is wrong only in the split.
\textbf{Verified:} closes to $4.8\times10^{-16}$ relative in
\code{tests/test\_annular.py} (\code{docs/OPEN\_QUESTIONS.md} Round 5).

\subsection{Rod Bundle Correction factors}

\code{correlations/bundle.py::Bundle}. Two interchangeable $\psi$ factors converting a
round-tube heat transfer coefficient to a rod-bundle one, $htc_{pin}=\psi\cdot
htc_{round\,tube}$ (Hughes et al.\ 2014, Eq.\ 11):

\textbf{Weissman:} $\psi = c_1 R + c_2$, $c_1=1.826$, $c_2=-1.0430$, $R=P/D$.
\textbf{Valid range / uncertainty / reference: \notest{}} --- named only as
``Weissman'' in the original source.

\textbf{Presser:} $\psi = c_1 + c_2 R - c_3 e^{-7(R-1)}$, $c_1=0.9217$, $c_2=0.1478$,
$c_3=0.1130$. \textbf{Reference:} Presser (1967), via Hughes et al.\ (2014) Eq.\ (10).
\textbf{Valid range:} \notest{} numerically --- exercised by Hughes at $P/D=1.15$
($\psi=0.94$) and $1.28$ ($\psi=0.98$); $\psi$ is \emph{not} monotone near unity and
turns over into a penalty ($\psi<1$) below roughly $P/D=1.05$, which is not itself
confirmed as a physical effect or an artifact of the fit running out of data (see
\code{docs/OPEN\_QUESTIONS.md} Q44).

\textbf{Physics fix, applied in two places, one per commit:} Presser was not applied
anywhere before this cleanup, despite both single-channel solvers using rod-bundle
geometry on at least one coolant channel.
\begin{itemize}
  \item \code{sca/rod.py} (commit \code{8cc40d9}): the whole channel is a square-pitch
    rod bundle; Presser is now applied to its one convective coefficient.
  \item \code{sca/annular.py} (commit \code{76d823e}): applied to the \textbf{outer
    channel only} --- a square-pitch cell around the cladding OD, the bundle geometry
    Presser was fit for --- and \emph{not} the inner channel, which is a bored tube
    through the pellet, the geometry Swenson itself was fit on. Measured on the default
    geometry (5~kW/m, 20 nodes): $htc_{conv,o}$ max $20801.8\to20721.1$~W/m$^2$-K,
    $T_{fo,o}$ max $677.975\to677.988$~K, flux-split closure unchanged at
    $4.75\times10^{-16}$. The effect is small only because the default lattice
    ($P/D=1.048$) sits almost exactly where $\psi\approx1$ (0.996); at a more typical
    $P/D=1.3$ the correction is $+10\%$.
\end{itemize}

Neither Weissman nor Presser drifted between the pre-cleanup copies (no numeric defect
found, D8).

% =====================================================================================
\section{Monte Carlo Subroutine and Error Propagation}
% =====================================================================================

\code{pinthac/uncertainty.py}. One mechanism, applied uniformly to both the property
libraries and the correlation libraries, per the brief's ``keep the API dead simple''
instruction:

\begin{itemize}
  \item \code{enable(seed)} / \code{disable()}: a module-level on/off switch with a
    single seed, so a Monte Carlo study is reproducible without every caller seeding
    numpy and torch's global generators itself (which would also entangle the
    library's randomness with, e.g., a training loop's weight initialization). Off by
    default: a deterministic run must never silently receive a perturbed answer, and
    must not pay for random-number generation it did not ask for.
  \item \code{perturb(value, rel\_sigma)}: the single mechanism behind every
    uncertainty band this library produces,
    \[
      value_{perturbed} = value\cdot(1+rel\_sigma\cdot z), \qquad z\sim N(0,1)
    \]
    multiplicative (not additive) so a positive quantity stays positive and because
    correlation uncertainties are published as percentages. One independent $z$ per
    element, so a batched call is equivalent to many independent Monte Carlo trials at
    once. The draw is treated as a constant with respect to autograd --- a sample of
    modelling error, not a differentiable function of the inputs --- so a gradient
    through a perturbed result is the underlying model's gradient scaled by that trial's
    realized factor, the correct sensitivity for that trial.
  \item \code{band(value, rel\_sigma, n\_samples)}: draws \code{n\_samples} independent
    perturbations of one value, stacked along a new leading axis, for the plotting path
    (a figure showing an uncertainty band needs many samples of the same curve, which
    \code{perturb} alone --- one perturbation per element --- cannot give). Works
    regardless of the module switch, since requesting a band is itself an explicit
    request for perturbed values.
\end{itemize}

\textbf{Consumer:} \code{figures/liquid\_metal\_uncertainty.py} --- thermal conductivity
vs.\ temperature for sodium, lead and LBE, with a 95\% Monte Carlo band (4000 trials per
temperature) built from each metal's own \code{uncert\_k} (linearly interpolated between
its stated low/high bound across its valid range, a stated modeling choice, not a value
given directly by Sobolev). Range checking that feeds the same uncertainty tables lives
in \code{pinthac/ranges.py::check(model\_name, values, table)}: a shared
\code{warnings.warn} helper, never an exception, raised at most once per call regardless
of batch size.

% =====================================================================================
\section*{What this manual deliberately does not cover}
% =====================================================================================

Per \code{docs/DECISIONS.md} (``Scope: transient is elsewhere''), transient
single-channel analysis (both a finite-difference solver and a physics-informed neural
network over time) lives in a separate repository and is out of scope here; nothing in
\code{pinthac/sca/} has a time derivative. Every neural surrogate documented above
(\code{ml/pinn.py}, \code{ml/deeponet.py}) is steady-state, parametrized over axial
position or operating conditions, not time. Two-phase (PWR/BWR) single-channel
bookkeeping --- onset-of-nucleate-boiling detection, quality/void-fraction tracking,
subcooled-boiling closure --- is likewise not implemented in either
\code{sca/rod.py} or \code{sca/annular.py}, even though the two-phase heat transfer
\emph{correlations} above (Schrock-Grossman, Chen, Bjorge) are implemented and
selectable by name.

\end{document}
